\subsubsection{Abel 有限部常数 logM 的接口式}\label{app:real-input-40-finite-part}
令 primitive 轨道生成函数
\[
P(z):=\sum_{n\ge 1}p_n z^n,
\qquad
\mathcal{P}(r):=P(r/\lambda)=\sum_{n\ge 1}\frac{p_n}{\lambda^n}r^n.
\]
则在 $r\uparrow 1$ 时有标准奇性展开
\[
\mathcal{P}(r)=-\log(1-r)+\log\mathfrak{M}+o(1),
\]
从而 Abel 有限部可写为
\[
\log\mathfrak{M}=\mathrm{fp}_{r=1}\bigl(\mathcal{P}(r)+\log(1-r)\bigr),
\]
这与主文第 \ref{sec:zeta-finite-part} 节的 Abel/有限部接口完全同构。

\begin{remark}[prime-orbit Mertens Abel 和（定义 \ref{def:abel-finite-part} 的同口径实例）]\label{rem:real-input-40-mertens-abel}
在本核上，$\mathcal{P}(r)$ 亦可直接视为 prime-orbit Mertens 和的 Abel 生成函数。若定义
\[
\mathcal{M}_{\mathrm{Abel}}(r):=\sum_{n\ge 1}\frac{p_n}{\lambda^n}r^n,
\]
则 $\mathcal{M}_{\mathrm{Abel}}(r)=\mathcal{P}(r)$，并且
\[
\mathcal{M}_{\mathrm{Abel}}(r)=-\log(1-r)+\log\mathfrak{M}+o(1)\qquad(r\uparrow 1).
\]
因此
\[
\boxed{\ \log\mathfrak{M}=\mathrm{fp}_{r=1}\Bigl(\mathcal{M}_{\mathrm{Abel}}(r)+\log(1-r)\Bigr)\ },
\]
这正是主文定义 \ref{def:abel-finite-part}“Abel 正则化 + Laurent 常数项有限部”的同型读出：区别仅在于此处的临界 $1/n$ 发散来自 primitive 轨道权 $p_n/\lambda^n\sim 1/n$，而非轨道 Dirichlet 级数的 $n^{-s}$ 临界。
\end{remark}

\begin{corollary}[prime-orbit Mertens 常数与 Abel 有限部常数的恒等式]\label{cor:real-input-40-Mertens-constant}
\leanverified{paper\_real\_input\_40\_mertens\_constant}
令
\[
\mathcal{M}(N):=\sum_{n\le N}\frac{p_n}{\lambda^n}.
\]
则存在常数 $\mathsf{M}$ 使得
\[
\mathcal{M}(N)=\log N+\mathsf{M}+o(1),
\qquad
\boxed{\ \mathsf{M}=\gamma+\log\mathfrak{M}\ }.
\]
\end{corollary}

\begin{proof}
写
\(
\frac{p_n}{\lambda^n}=\frac{1}{n}+b_n
\)
其中 $b_n:=\frac{p_n}{\lambda^n}-\frac{1}{n}$。
由 $p_n=\lambda^n/n+O(\varphi^n/n)$（见上文 $|\mathrm{spec}_2|=\varphi<\lambda$）可知 $\sum_{n\ge 1}b_n$ 绝对收敛。
于是
\[
\sum_{n\le N}\frac{p_n}{\lambda^n}
\;=\;
\sum_{n\le N}\frac{1}{n}
\;+\;
\sum_{n\le N}b_n
\;=\;
\log N+\gamma+\sum_{n\ge 1}b_n+o(1).
\]
另一方面，由备注 \ref{rem:real-input-40-mertens-abel} 中的 Abel 有限部读出，
\(
\sum_{n\ge 1}b_n=\log\mathfrak{M}
\)
（等价于 $\mathrm{fp}_{r=1}\bigl(\sum b_n r^n\bigr)=\sum b_n$），从而得到 $\mathsf{M}=\gamma+\log\mathfrak{M}$。
\end{proof}

\paragraph{Möbius--Abel 快速公式（接口计算式）}
把 Möbius--log 反演写到 Abel 坐标中，可得
\[
\log\mathfrak{M}
:=\log C+\sum_{m\ge 2}\frac{\mu(m)}{m}\log\widetilde{\zeta}_M(\lambda^{1-m}),
\qquad (\lambda^{1-m}\in(0,1)\ \text{且指数趋零}).
\]
该表示的每一项仅需要对 $\widetilde{\zeta}_M$ 在收敛圆盘内部的采样，因而数值收敛极快。
等价地，由 $z_\star=\lambda^{-1}$ 与 $\widetilde{\zeta}_M(\lambda^{1-m})=\zeta_M(z_\star^{\,m})=\zeta_M(\lambda^{-m})$，亦可写为
\[
\log\mathfrak{M}
:=\log C+\sum_{m\ge 2}\frac{\mu(m)}{m}\log\zeta_M(\lambda^{-m}).
\]

\begin{proposition}[Möbius--Abel 快速公式的显式尾项证书：有限部常数可指数精度截断]\label{prop:real-input-40-logM-trunc-bound}
设 $\zeta_M(z)=1/\det(I-zM)$ 在圆盘 $|z|\le z_\star^{\,2}=\lambda^{-2}$ 内解析且 $\zeta_M(0)=1$。定义常数
\[
A:=\sup_{0<|z|\le \lambda^{-2}}\ \frac{|\log\zeta_M(z)|}{|z|}\ <\ \infty.
\]
对任意截断阶 $K\ge 2$，令
\[
\log\mathfrak{M}^{(K)}:=\log C+\sum_{m=2}^{K}\frac{\mu(m)}{m}\log\zeta_M(\lambda^{-m}).
\]
则尾项满足显式误差上界
\[
\boxed{\;
|\log\mathfrak{M}-\log\mathfrak{M}^{(K)}|
\le
A\sum_{m>K}\frac{\lambda^{-m}}{m}
\le
\frac{A}{(K+1)(\lambda-1)}\,\lambda^{-K}.
\;}
\]
\end{proposition}
\begin{proof}
对 $m\ge 2$，点 $\lambda^{-m}$ 落在圆盘 $|z|\le\lambda^{-2}$ 内，故由 $A$ 的定义得
\(
|\log\zeta_M(\lambda^{-m})|\le A\lambda^{-m}
\)。
又 $|\mu(m)|\le 1$，因此
\[
|\log\mathfrak{M}-\log\mathfrak{M}^{(K)}|
\le \sum_{m>K}\frac{1}{m}|\log\zeta_M(\lambda^{-m})|
\le A\sum_{m>K}\frac{\lambda^{-m}}{m}.
\]
对尾和用单调性估计：
\[
\sum_{m>K}\frac{\lambda^{-m}}{m}
\le
\frac{1}{K+1}\sum_{m>K}\lambda^{-m}
:=\frac{1}{K+1}\cdot\frac{\lambda^{-(K+1)}}{1-\lambda^{-1}}
:=\frac{1}{(K+1)(\lambda-1)}\,\lambda^{-K}.
\]
证毕。
\end{proof}
\leanverified{paper\_real\_input\_40\_logM\_trunc\_bound}

\begin{remark}[对数精度截断深度：常数层可在 \texorpdfstring{$O(\log(1/\varepsilon))$}{O(log(1/eps))} 次采样内证书化]\label{rem:real-input-40-logM-logeps-truncation}
由命题 \ref{prop:real-input-40-logM-trunc-bound} 的尾项界，给定 \(\varepsilon\in(0,1)\)，取 \(K\ge 2\) 使得
\[
\frac{A}{(K+1)(\lambda-1)}\,\lambda^{-K}\le \varepsilon
\]
即可保证 \(|\log\mathfrak{M}-\log\mathfrak{M}^{(K)}|\le \varepsilon\)。
因此在固定 \(\lambda>1\) 的设定下，达到对数精度 \(\varepsilon\) 所需的采样次数满足 \(K\asymp \log(1/\varepsilon)\)，并且每一项仅涉及收敛圆盘内部点 \(\lambda^{-m}\) 上的 \(\log\zeta_M\) 采样。
\end{remark}

\begin{proposition}[有限部常数的稳定性泛函：采样误差仅以 Möbius 权线性聚合]\label{prop:real-input-40-logM-stability-functional}
沿用命题 \ref{prop:real-input-40-logM-trunc-bound} 的设定。若对每个 \(m\ge 2\) 的采样存在加性误差 \(\epsilon_m\)，使
\[
\log\zeta_M(\lambda^{-m})\ \longmapsto\ \log\zeta_M(\lambda^{-m})+\epsilon_m,
\]
则有限部常数的漂移满足严格上界
\[
\boxed{\;
|\Delta\log\mathfrak{M}|\le \sum_{m\ge 2}\frac{|\epsilon_m|}{m}\; }.
\]
特别地，若存在 \(\varepsilon>0\) 使对一切 \(m\ge 2\) 有 \(|\epsilon_m|\le \varepsilon\,\lambda^{-m}\)，则
\[
|\Delta\log\mathfrak{M}|
\le
\varepsilon\sum_{m\ge 2}\frac{\lambda^{-m}}{m}
\le
\frac{\varepsilon\,\lambda^{-2}}{2(1-\lambda^{-1})}.
\]
\leanverified{paper\_real\_input\_40\_logM\_stability\_functional}
\end{proposition}
\begin{proof}
由 Möbius--Abel 表示 \(\log\mathfrak{M}=\log C+\sum_{m\ge 2}\frac{\mu(m)}{m}\log\zeta_M(\lambda^{-m})\) 得
\(
\Delta\log\mathfrak{M}=\sum_{m\ge 2}\frac{\mu(m)}{m}\epsilon_m
\)；
结合 \(|\mu(m)|\le 1\) 取绝对值即得第一式。
若 \(|\epsilon_m|\le \varepsilon\lambda^{-m}\)，则
\[
|\Delta\log\mathfrak{M}|
\le \varepsilon\sum_{m\ge 2}\frac{\lambda^{-m}}{m}
\le \frac{\varepsilon}{2}\sum_{m\ge 2}\lambda^{-m}
=\frac{\varepsilon\,\lambda^{-2}}{2(1-\lambda^{-1})}.
\]
证毕。
\end{proof}

\paragraph{与 primitive 分解对齐的双重 Möbius--Abel 级数}
由分解恒等式
\[
\log\zeta_M(z)=\sum_{m\ge 1}\frac{P(z^m)}{m}
\]
并将 $z=r/\lambda$ 代入、把 $m=1$ 项单独提出，可得
\[
P(r/\lambda)=-\log(1-r)+\log\mathfrak{M}+o(1)\qquad(r\uparrow 1),
\]
以及常数分解
\[
\log\mathfrak{M}=\log C-\sum_{m\ge 2}\frac{P(\lambda^{-m})}{m}.
\]
另一方面，Möbius 反演给出 primitive 生成函数的闭式
\[
P(z)=\sum_{k\ge 1}\frac{\mu(k)}{k}\log\zeta_M(z^k).
\]
将其代回即得到完全由 $\zeta_M$ 构成、且收敛极快的双和表示：
\[
\log\mathfrak{M}
:=\log C-\sum_{m\ge 2}\frac{1}{m}\sum_{k\ge 1}\frac{\mu(k)}{k}\log\zeta_M(\lambda^{-mk}).
\]

\paragraph{\texorpdfstring{$\mu$}{mu}-Pochhammer 分解：把 primitive 复杂度压缩为若干通用乘积函数}
Möbius--$\log$ 反演本质上把一个有限维有理 \(\zeta_M\) 转化为 primitive 生成函数 \(P(z)\) 的对数型组合；为显式化其“复杂性来源”，引入一族带 Möbius 权的 Pochhammer 乘积。

\begin{definition}[\texorpdfstring{$\mu$}{mu}-Pochhammer 乘积]\label{def:mu-pochhammer}
对任意复参数 \(\alpha\in\CC\)，在 \(|z|<1\) 的主值对数分支下定义
\[
\Phi_\alpha(z)
:\!=\prod_{m\ge 1}\bigl(1-\alpha z^m\bigr)^{\mu(m)/m},
\qquad
\log\Phi_\alpha(z)
:\!=\sum_{m\ge 1}\frac{\mu(m)}{m}\log\bigl(1-\alpha z^m\bigr).
\]
\end{definition}

\begin{proposition}[\texorpdfstring{$\mu$}{mu}-Pochhammer 分解公式]\label{prop:mu-pochhammer-decomposition}
\leanverified{paper\_mu\_pochhammer\_decomposition}
设 \(M\) 为有限维矩阵，\(\zeta_M(z)=1/\det(I-zM)\)，其非零特征值（按代数重数计）为 \(\lambda\neq 0\)，并记每个 \(\lambda\) 的代数重数为 \(m_\lambda\)。
则在 \(|z|\) 充分小（从而 \(\log\) 取主值且级数一致收敛）时有严格恒等式
\[
\boxed{\;
P(z)
=\sum_{k\ge 1}\frac{\mu(k)}{k}\log\zeta_M(z^k)
=-\sum_{\lambda\neq 0}m_\lambda\,\log\Phi_\lambda(z)\; }.
\]
\end{proposition}
\begin{proof}
由谱分解（按代数重数）有
\(
\det(I-zM)=\prod_{\lambda\neq 0}(1-\lambda z)^{m_\lambda}
\)，
因此
\(
\log\zeta_M(z)
=-\sum_{\lambda\neq 0}m_\lambda\log(1-\lambda z)
\)。
将其代入 Möbius--$\log$ 反演
\(
P(z)=\sum_{k\ge 1}\mu(k)\log\zeta_M(z^k)/k
\)，并交换有限和与一致收敛级数的求和次序，得到
\[
P(z)
=-\sum_{\lambda\neq 0}m_\lambda\sum_{k\ge 1}\frac{\mu(k)}{k}\log(1-\lambda z^k)
=-\sum_{\lambda\neq 0}m_\lambda\,\log\Phi_\lambda(z).
\]
证毕。
\end{proof}

\begin{proposition}[\texorpdfstring{$\mu$}{mu}-Pochhammer 的 necklace 展开]\label{prop:mu-pochhammer-necklace-expansion}
\leanverified{paper\_mu\_pochhammer\_necklace\_expansion}
在定义 \ref{def:mu-pochhammer} 的设定下，对 \(|z|<1\) 的主值对数分支，有严格幂级数展开
\[
\boxed{\;
\log\Phi_\alpha(z)
=-\sum_{n\ge 1}N_n(\alpha)\,z^n,
\qquad
N_n(\alpha):=\frac{1}{n}\sum_{d\mid n}\mu(d)\,\alpha^{n/d}\; }.
\]
\end{proposition}

\begin{proof}
由定义 \ref{def:mu-pochhammer}，
\(
\log\Phi_\alpha(z)=\sum_{m\ge 1}\frac{\mu(m)}{m}\log(1-\alpha z^m)
\)。
当 \(|z|<1\) 时，
\(
\log(1-\alpha z^m)=-\sum_{k\ge 1}\frac{\alpha^k}{k}z^{mk}
\)
在任意紧子圆盘上一致收敛；交换求和并按 \(n=mk\) 分组得
\[
\log\Phi_\alpha(z)
=-\sum_{n\ge 1}\frac{1}{n}\left(\sum_{m\mid n}\mu(m)\alpha^{n/m}\right)z^n,
\]
即得所示公式。证毕。
\end{proof}

\begin{corollary}[\texorpdfstring{$\Phi_1(z)=e^{-z}$}{Phi1=e(-z)} 的退化律]\label{cor:mu-pochhammer-phi1-exp-minus-z}
当 \(|z|<1\) 时有严格恒等式 \(\Phi_1(z)=e^{-z}\)。
\end{corollary}

\begin{proof}
由命题 \ref{prop:mu-pochhammer-necklace-expansion}，当 \(\alpha=1\) 时
\(
N_n(1)=\frac{1}{n}\sum_{d\mid n}\mu(d)
\)
仅在 \(n=1\) 时为 \(1\)，在 \(n\ge 2\) 时为 \(0\)。
因此 \(\log\Phi_1(z)=-z\)，即 \(\Phi_1(z)=e^{-z}\)。证毕。
\end{proof}

\begin{proposition}[necklace 系数的 Dirichlet--polylog/\texorpdfstring{$\zeta$}{zeta} 因子分解]\label{prop:mu-pochhammer-necklace-dirichlet-polylog}
\leanverified{paper\_mu\_pochhammer\_necklace\_dirichlet\_polylog}
令 \(N_n(\alpha)\) 如命题 \ref{prop:mu-pochhammer-necklace-expansion}。
当 \(|\alpha|<1\) 且 \(\Re(s)>0\) 时，有 Dirichlet 级数恒等式
\[
\boxed{\;
\sum_{n\ge 1}\frac{N_n(\alpha)}{n^s}
=\frac{\mathrm{Li}_{s+1}(\alpha)}{\zeta(s+1)}\; }.
\]
\end{proposition}

\begin{proof}
设 \(g(n):=\sum_{d\mid n}\mu(d)\alpha^{n/d}\)，则 \(N_n(\alpha)=g(n)/n\)。
当 \(|\alpha|<1\) 且 \(\Re(w)>1\) 时，绝对收敛给出
\[
\sum_{n\ge 1}\frac{g(n)}{n^w}
=\sum_{n\ge 1}\sum_{d\mid n}\frac{\mu(d)\alpha^{n/d}}{n^w}
=\sum_{d\ge 1}\frac{\mu(d)}{d^w}\sum_{m\ge 1}\frac{\alpha^m}{m^w}
=\frac{1}{\zeta(w)}\,\mathrm{Li}_{w}(\alpha),
\]
其中 \(\mathrm{Li}_{w}(\alpha)=\sum_{m\ge 1}\alpha^m/m^w\)。
令 \(w=s+1\) 并注意
\(
\sum_{n\ge 1}N_n(\alpha)/n^s=\sum_{n\ge 1}g(n)/n^{s+1}
\)，即得结论。证毕。
\end{proof}

\begin{theorem}[Abel 有限部常数的 \texorpdfstring{$\mu$}{mu}-Pochhammer 谱分解闭式]\label{thm:finite-part-mu-pochhammer-spectral-closed-form}
设 \(M\) 为有限维矩阵，\(\zeta_M(z)=\det(I-zM)^{-1}\)，并令 \(\lambda=\rho(M)>1\) 为 Perron 根。
假设 \(\lambda\) 为简单特征值且满足谱隙 \(|\lambda_j|<\lambda\)（对所有非零特征值 \(\lambda_j\neq\lambda\)，按代数重数计）。
令
\[
P(z):=\sum_{k\ge 1}\frac{\mu(k)}{k}\log\zeta_M(z^k),
\qquad
\mathcal{P}(r):=P(r/\lambda).
\]
并设 Abel 有限部常数 \(\log\mathfrak{M}\) 由
\[
\mathcal{P}(r)=-\log(1-r)+\log\mathfrak{M}+o(1)\qquad(r\uparrow 1)
\]
所定义。
再定义去奇异 \texorpdfstring{$\mu$}{mu}-Pochhammer
\[
\widetilde{\Phi}_\alpha(z):=\frac{\Phi_\alpha(z)}{1-\alpha z}
=\prod_{m\ge 2}\bigl(1-\alpha z^m\bigr)^{\mu(m)/m}.
\]
则有显式闭式
\[
\boxed{\;
\log\mathfrak{M}
=-\log\widetilde{\Phi}_{\lambda}\!\Bigl(\frac{1}{\lambda}\Bigr)
\;-\!\!\sum_{\substack{\lambda_j\neq 0\\ \lambda_j\neq\lambda}}
m_j\,\log\Phi_{\lambda_j}\!\Bigl(\frac{1}{\lambda}\Bigr)\; } ,
\]
其中 \(\{\lambda_j\neq 0\}\) 为 \(M\) 的非零特征值多重集，\(m_j\) 为代数重数。
\end{theorem}
\leanverified{paper\_finite\_part\_mu\_pochhammer\_spectral\_closed\_form}

\begin{proof}
由命题 \ref{prop:mu-pochhammer-decomposition} 与 \(\lambda\) 的简单性，有
\[
P(z)=-\log\Phi_{\lambda}(z)-\!\!\sum_{\substack{\lambda_j\neq 0\\ \lambda_j\neq\lambda}}m_j\log\Phi_{\lambda_j}(z).
\]
当 \(z\to 1/\lambda\) 时，对任意 \(\lambda_j\neq\lambda\) 与任意 \(m\ge 1\) 有
\(
|\lambda_j z^m|
\le |\lambda_j|\lambda^{-m}
\le (|\lambda_j|/\lambda)\lambda^{1-m}<1
\)，
故 \(\log\Phi_{\lambda_j}(z)\) 在 \(z\to 1/\lambda\) 时收敛到 \(\log\Phi_{\lambda_j}(1/\lambda)\)。
\par
对 Perron 根项，分离 \(m=1\) 因子得
\[
\log\Phi_{\lambda}(z)=\log(1-\lambda z)+\log\widetilde{\Phi}_{\lambda}(z),
\]
并且 \(\log\widetilde{\Phi}_{\lambda}(z)\) 在 \(z\to 1/\lambda\) 时收敛到 \(\log\widetilde{\Phi}_{\lambda}(1/\lambda)\)。
取 \(z=r/\lambda\)（\(r\uparrow 1\)）并代回，得到
\[
\mathcal{P}(r)
=-\log(1-r)
\;-\log\widetilde{\Phi}_{\lambda}(1/\lambda)
\;-\!\!\sum_{\substack{\lambda_j\neq 0\\ \lambda_j\neq\lambda}}m_j\log\Phi_{\lambda_j}(1/\lambda)
o(1).
\]
与 \(\mathcal{P}(r)=-\log(1-r)+\log\mathfrak{M}+o(1)\) 比较即得结论。证毕。
\end{proof}


\begin{corollary}[真实输入 \(40\) 状态核的显式分解]\label{cor:real-input-40-mu-pochhammer-explicit}
\leanverified{paper\_real\_input\_40\_mu\_pochhammer\_explicit}
在命题 \ref{prop:real-input-40-essential-20} 的真实输入 \(40\) 状态核上，非零特征值（含重数）为
\[
\{\varphi^2,-\varphi,1^{(2)},(-1)^{(3)},\varphi^{-2},\varphi^{-1}\}.
\]
因此 primitive 生成函数满足完全显式的分解
\[
\boxed{
P(z)=
-\log\Phi_{\varphi^2}(z)
-\log\Phi_{-\varphi}(z)
-2\log\Phi_{1}(z)
-3\log\Phi_{-1}(z)
-\log\Phi_{\varphi^{-2}}(z)
-\log\Phi_{\varphi^{-1}}(z).
}
\]
进一步由引理 \ref{lem:mobius-collapse} 的闭式 \(\log\Phi_{1}(z)=-z\)、\(\log\Phi_{-1}(z)=z-z^2\)，可将“平凡谱”部分精确压缩为二次多项式：
\[
\boxed{
P(z)=
-\log\Phi_{\varphi^2}(z)
-\log\Phi_{-\varphi}(z)
-\log\Phi_{\varphi^{-2}}(z)
-\log\Phi_{\varphi^{-1}}(z)
-z+3z^2.
}
\]
\end{corollary}

\paragraph{自然边界：有理 \texorpdfstring{$\zeta$}{zeta} 与“奇点森林”的机制链}
上式把 \(P\) 的解析复杂度集中到 \(|\lambda|>1\) 的两项 \(\Phi_{\varphi^2}\) 与 \(\Phi_{-\varphi}\) 上：它们把 \(\zeta_M\) 的有限多主极点，经 Möbius 权的根系复制，生成一族在单位圆上稠密聚点的对数分支奇点。

\begin{theorem}[单位圆自然边界]\label{thm:primitive-natural-boundary-unit-circle}
设 \(P(z)\) 为真实输入 \(40\) 状态核的 primitive 生成函数。则 \(P\) 从 \(|z|<1\) 的解析延拓在 \(|z|=1\) 处存在自然边界：
任取单位圆上一点 \(e^{i\theta}\) 及其任意邻域，都包含 \(P\) 的对数分支奇点；因此 \(P\) 不可能解析延拓穿过 \(|z|=1\) 的任何非平凡弧段。
\leanverified{paper\_primitive\_natural\_boundary\_unit\_circle}
\end{theorem}
\begin{proof}
由推论 \ref{cor:real-input-40-mu-pochhammer-explicit}，\(P\) 是若干项 \(-\log\Phi_\alpha(z)\) 与多项式之和，其中 \(\alpha\in\{\varphi^2,-\varphi,\varphi^{-2},\varphi^{-1}\}\)。
当 \(|\alpha|<1\) 时，\((1-\alpha z^m)\neq 0\) 对所有 \(|z|\le 1\) 成立，故 \(\log\Phi_\alpha\) 在 \(|z|\le 1\) 解析，从而不会在单位圆制造奇点。

当 \(|\alpha|>1\) 时（此处 \(\alpha=\varphi^2\) 或 \(\alpha=-\varphi\)），对任意平方自由整数 \(m\ge 2\) 与任意 \(m\) 次单位根 \(\omega\)，点
\[
z_{m,\omega}^{(\alpha)}:=\alpha^{-1/m}\omega
\]
满足 \(1-\alpha z^m=0\) 且 \(|z_{m,\omega}^{(\alpha)}|=|\alpha|^{-1/m}\uparrow 1\)（\(m\to\infty\)）。
由定义 \ref{def:mu-pochhammer}，\(\log\Phi_\alpha\) 的求和项中包含
\(\frac{\mu(m)}{m}\log(1-\alpha z^m)\)，而 \(\mu(m)\neq 0\) 当且仅当 \(m\) 平方自由。
因此在每个平方自由 \(m\) 上，\(\log\Phi_\alpha\) 在 \(z=z_{m,\omega}^{(\alpha)}\) 处具有不可去的对数分支奇点；从而 \(P\) 在这些点同样具有分支奇点（多项式部分解析且不会抵消分支）。

最后注意到：对固定 \(m\)，集合 \(\{\omega:\omega^m=1\}\) 在单位圆上等分；当 \(m\) 取遍平方自由整数时，这些等分点的并集在单位圆上稠密。
同时半径 \(|\alpha|^{-1/m}\) 取值趋于 \(1\)，故上述分支奇点在单位圆上处处有聚点，从而 \(|z|=1\) 为 \(P\) 的自然边界。证毕。
\end{proof}

\begin{remark}[单点单圈的算术单值差（monodromy）]\label{rem:primitive-monodromy-mu}
沿上述记号，绕 \(z_{m,\omega}^{(\varphi^2)}=\varphi^{-2/m}\omega\) 做一小圈解析延拓，
来自项 \(-\frac{\mu(m)}{m}\log(1-\varphi^2 z^m)\) 的分支切换给出
\[
\Delta P=-\frac{\mu(m)}{m}\cdot 2\pi i,
\]
其余各项在该邻域解析不变；绕 \(z_{m,\omega}^{(-\varphi)}=\varphi^{-1/m}e^{(2j+1)\pi i/m}\) 的情形同理。
因此 \(P\) 的分支结构在单位圆附近呈现显式的“平方自由根系晶格”，生成元权重直接由 \(\mu(m)/m\) 给出。
\end{remark}

\paragraph{Abel 常数的结构分解：PF 通用常数 \texorpdfstring{$+$}{+} 稳定因子 \texorpdfstring{$+$}{+} 代数修正项}
令 \(\lambda=\varphi^2\)、\(z_\star=\lambda^{-1}=\varphi^{-2}\)。由 \(P(z)=-\log(1-\lambda z)+\log\mathfrak{M}+o(1)\)（\(z\to z_\star^{-}\)）与推论 \ref{cor:real-input-40-mu-pochhammer-explicit}，可把 \(\log\mathfrak M\) 拆解为一项只依赖 PF 根的通用常数、三个快速收敛的稳定 \(\mu\)-Pochhammer 值，以及一个纯代数修正项。

\begin{corollary}[Abel 有限部常数的 \texorpdfstring{$\mu$}{mu}-Pochhammer 分解]\label{cor:logM-mu-pochhammer-split}
定义 PF 通用常数
\[
K_{\mathrm{PF}}
:=\prod_{m\ge 2}\bigl(1-\varphi^{2(1-m)}\bigr)^{\mu(m)/m},
\qquad
\log K_{\mathrm{PF}}
=\sum_{m\ge 2}\frac{\mu(m)}{m}\log\bigl(1-\varphi^{2(1-m)}\bigr).
\]
则在真实输入 \(40\) 状态核上有严格恒等式
\[
\boxed{
\log\mathfrak{M}
=-\log K_{\mathrm{PF}}
-\log\Phi_{-\varphi}(z_\star)
-\log\Phi_{\varphi^{-1}}(z_\star)
-\log\Phi_{\varphi^{-2}}(z_\star)
\;+\;\bigl(-z_\star+3z_\star^2\bigr).
}
\]
并且代数修正项可完全化简为
\[
-z_\star+3z_\star^2=-\varphi^{-2}+3\varphi^{-4}=9-4\sqrt5.
\]
\leanverified{paper\_logm\_mu\_pochhammer\_split}
\end{corollary}
